% Chapter - Stochasticity and Markov Processes ........................................................................
\chapter{Stochasticity and Markov processes}
\label{chapter7}

\epigraph{\textit{The theory of probabilities is at bottom nothing\\but common sense reduced to calculation.}}{— Pierre-Simon Laplace}

There is a certain comfort in the assumption of independence. If the outcome of each experiment tells us nothing about the next, then the world becomes, in a sense, memoryless: the past dissolves, and only the present matters. Much of classical probability and statistics, as we have developed it in previous chapters, rests on this assumption. Samples are drawn independently, errors are uncorrelated, observations do not influence one another. These assumptions are not merely mathematical conveniences---they reflect a deeper philosophical position about how randomness operates in the world.

\medskip

The mathematization of uncertainty began precisely with this idealized picture. Bernoulli's law of large numbers established that stable regularities emerge from repeated random trials, revealing that randomness itself could obey precise quantitative laws and need not imply disorder \cite{bernoulli_1713_ars}. Laplace extended this insight beyond gambling to scientific reasoning more broadly, arguing that probability expresses rational inference under incomplete knowledge, and that apparent randomness often reflects limited information rather than metaphysical indeterminacy---thereby embedding uncertainty within a wider epistemological framework \cite{laplace_1812_theorie}. On this view, the world is not fundamentally chaotic; it is merely incompletely known, and probability is the language of that incompleteness.

\medskip

Yet a moment's reflection reveals how rarely the assumption of independence holds in practice. The weather tomorrow is not independent of the weather today. The price of a stock this afternoon is shaped by its price this morning. The word that follows in a sentence is constrained by the words that came before. The health of a patient next week depends on their condition this week. In all these cases, the present state of a system carries information about its future evolution, and ignoring this dependence would lead to poor models and worse predictions.

\medskip

During the nineteenth and early twentieth centuries, attention shifted from isolated probabilistic events to sequences of dependent variables evolving over time. The question was no longer merely what the probability of a single event is, but how a system changes, accumulates, and forgets as it moves through a succession of states. It is in this context that one of the more unusual disputes in the history of mathematics unfolded---one that was simultaneously about probability theory, theology, and the nature of human freedom.

\medskip

Pavel Nekrasov was a mathematician, theologian, and conservative intellectual who held a privileged position in the Russian academic establishment of the late nineteenth century, eventually becoming rector of Moscow University and a close ally of tsarist institutions. His central mathematical claim was that the law of large numbers required the events in question to be \textit{independent}. Nekrasov did not advance this claim merely as a technical observation. He argued that it had profound metaphysical consequences: if stable social regularities such as crime rates or patterns of collective behavior obeyed probabilistic laws, and if such laws required independence, then the individuals comprising society must each be making genuinely free and unconstrained choices. The law of large numbers, on this reading, was not merely a theorem about coin tosses---it was mathematical proof of free will, and by extension, of the moral foundations of Christian society and the legitimacy of the tsarist order. Mathematics, in Nekrasov's hands, had become a weapon of conservative theology.

\medskip

Andrey Andreyevich Markov (1856--1922) found this argument deeply objectionable, on both mathematical and personal grounds. Markov was, by all accounts, a man of fierce temperament and uncompromising convictions---an atheist, a socialist sympathizer, and an outspoken critic of the Orthodox Church and the tsarist establishment. When Lev Tolstoy was excommunicated from the Russian Orthodox Church in 1901, Markov publicly demanded his own excommunication in protest, a gesture of defiant solidarity that caused a considerable scandal. He was not a man inclined to let a mathematical error stand, least of all one dressed in theological clothing and deployed in support of institutions he despised. Nekrasov's claim that independence was necessary for the law of large numbers was, in Markov's view, simply wrong---and he set about proving it \cite{markov_1906_extension}.

\medskip

Markov's strategy was to construct an explicit counterexample: a sequence of \textit{dependent} events that nonetheless exhibited the stable long-run regularities guaranteed by the law of large numbers. Rather than work with an artificial mathematical construction, he chose a literary one. He took the opening chapters of Pushkin's celebrated novel in verse \textit{Eugene Onegin} and proceeded to analyze the text with meticulous care, manually counting, letter by letter, the transitions between vowels and consonants across thousands of characters. For each letter, he recorded whether it was a vowel (V) or a consonant (C), and then tallied how often a vowel was followed by another vowel, how often by a consonant, how often a consonant was followed by a vowel, and so on. What emerged was a table of transition counts---the raw material for what we would now call a two-state Markov chain, with states V and C and transition probabilities estimated from the data.

\medskip

The point was precise and devastating. The letters of a text are manifestly not independent: knowing that the current letter is a vowel makes it more or less likely that the next is a consonant, depending on the phonological structure of the language. Dependence is everywhere. And yet the long-run proportions of vowels and consonants stabilize perfectly well, just as the law of large numbers would predict for independent events. Independence, Markov showed, was not a necessary condition for probabilistic regularity---it was merely a sufficient one. Nekrasov's theological argument collapsed with his mathematics.

\medskip

What Markov had done, almost as a polemical aside, was to inaugurate the systematic study of dependent stochastic sequences. The abstraction is straightforward once the example is in hand: instead of tracking the actual letters of a text, one retains only their classification into states---V or C---and studies how the system moves between these states over time. Each state carries a probability of transitioning to each other state, and the sequence of states forms what we now call a \textit{Markov chain}. The specific literary content of \textit{Eugene Onegin} drops away entirely; what remains is a mathematical structure of extraordinary generality and power.

\medskip

The story does not end there. Several decades later, on the other side of the world, a different kind of mathematical problem was generating a different kind of urgency. The year was 1946, and Stanislaw Ulam was recovering from a serious illness at Los Alamos, passing the time by playing solitaire. He found himself wondering about a deceptively simple question: what is the probability that a game of solitaire comes out successfully? The combinatorics were intractable---the number of possible card arrangements is astronomical---but Ulam realized that one could simply \textit{play} the game many times and observe the fraction of successes. The answer could be estimated by simulation rather than derived by calculation. He brought the idea to John von Neumann, who immediately recognized its potential for the nuclear calculations then underway at Los Alamos, where the behavior of neutrons in fissile material presented exactly the same kind of high-dimensional intractability. The method needed a name, and Nicholas Metropolis provided one: Monte Carlo, after the famous casino in Monaco, with a nod to Ulam's uncle who was known to frequent it \cite{ulam_1976_adventures, metropolis_1949_monte}.

\medskip

What connected Ulam's solitaire game to Markov's vowel counts was not immediately obvious, but the connection runs deep. Both are fundamentally concerned with sequences that evolve step by step according to probabilistic rules. If the steps of that sequence are designed carefully---so that the long-run distribution of states matches a desired target---then simulating the sequence produces samples from that target distribution, even when the target is far too complex to work with analytically. This insight, that Markov chains could be used as engines of computation and not merely as objects of mathematical study, gave rise to the family of methods known as Markov chain Monte Carlo, which we will explore in detail later. It is one of the more remarkable convergences in the history of applied mathematics: a Russian atheist counting letters in a poem, and an ailing Polish mathematician playing cards in a New Mexico laboratory, had each contributed an essential half of one of the most powerful computational tools in modern science.

\medskip

This temporal generalization of probability was later placed on rigorous foundations through Kolmogorov's axiomatization and Doob's development of stochastic process theory, which formalized random evolution as measurable structure on a probability space and established the analytical framework that continues to underlie modern stochastic analysis \cite{kolmogorov_1933_grundbegriffe, doob_1953_stochastic}. Scientific models often begin with deterministic equations, yet many natural systems exhibit variability that cannot be captured by fixed trajectories alone. Stochastic modelling represents not chaos but structured uncertainty governed by probabilistic laws. Variability over time or across states may be intrinsic rather than merely observational, and it requires explicit probabilistic structure rather than deterministic approximation.

\medskip

The question, then, is not whether dependence and variability exist, but how to model them in a tractable and mathematically rigorous way. The most natural first step is to ask: what if the future depends on the present, but not on the more distant past? What if, given where we are now, where we came from is irrelevant? This simple and elegant idea---born from a dispute about free will, proved on a Russian poem, and later harnessed to simulate the behavior of neutrons---is the foundation of the theory of Markov models.

\medskip

In this chapter, we trace the development of this framework from its origins to its modern incarnations. We begin with the fundamental mathematical definition of a Markov chain and its key properties, illustrated through simple and familiar examples. We then extend the framework to continuous time, where states evolve not in discrete steps but along a continuous timeline. We introduce hidden Markov models, where the underlying state is not directly observable but must be inferred from indirect evidence. We discuss Markov chain Monte Carlo methods, which harness the structure of Markov chains to solve otherwise intractable problems in Bayesian computation. We then turn to Markov decision processes and the theory of reinforcement learning, where an agent must learn to act optimally in an uncertain environment. Finally, we trace the conceptual thread that connects these classical ideas to the modern architecture of large language models and transformers, and reflect on the limitations of the Markov assumption itself. Throughout, the emphasis will be on intuition and ideas rather than technical machinery. The mathematics will be introduced as needed, but always in service of understanding---what these models say, why they work, and where they fall short.

% Section 1 ...........................................................................................................
\section{Stochasticity and the Markov property}

Few ideas in mathematics are as powerful as a well-chosen simplification. The theory of Markov chains rests on one such simplification: that the future of a system, given its present state, is independent of its past. This is known as the \textit{Markov property}, and while it may seem like a strong assumption, it captures the essential structure of an enormous range of real-world processes. As we saw in the introduction, Markov arrived at this idea by studying sequences of letters in a poem. The abstraction he performed---retaining only the state of the system at each step and the probabilities of moving between states---turns out to be applicable far beyond literary analysis, and it is worth developing carefully.

\medskip

To make this precise, consider a sequence of random variables $X_0, X_1, X_2, \dots$, each taking values in some set $\mathcal{S}$ called the \textit{state space}. This sequence represents the evolution of a system over time: $X_n$ is the state of the system at time $n$. The sequence is called a \textit{Markov chain} if, for all times $n$ and all states $x_0, x_1, \dots, x_{n+1} \in \mathcal{S}$,
\begin{equation}
    \mathbb{P}(X_{n+1} = x_{n+1} \mid X_n = x_n, X_{n-1} = x_{n-1}, \dots, X_0 = x_0) = \mathbb{P}(X_{n+1} = x_{n+1} \mid X_n = x_n) \; .
    \label{eq:markov_property}
\end{equation}
In words: given the present state $X_n$, the probability of the next state $X_{n+1}$ does not depend on how the system arrived at $X_n$. The past, once the present is known, becomes irrelevant. This is the \textit{Markov property}, also sometimes called the \textit{memoryless property}.

\medskip

It is worth pausing to appreciate what this assumption does and does not say. It does not claim that the past has no influence on the future---it clearly does, through the current state. What it says is that all the relevant information about the past is \textit{summarized} in the present state. If we know where we are, knowing the path we took to get here adds nothing. In Markov's vowel-consonant chain, for instance, knowing that the current letter is a vowel tells us everything the model needs to know about what comes next---we do not need to remember whether the letter before that was also a vowel, or what the opening line of the poem said. This is a modeling assumption, not a law of nature, and its appropriateness depends entirely on how we define the state of the system. Chosen well, the state summarizes the past; chosen poorly, the Markov property fails and the model misleads.

\subsection{Transition probabilities and the transition matrix}

The essential ingredient of a Markov chain is the collection of probabilities governing how the system moves from one state to another. For a finite state space $\mathcal{S} = \{1, 2, \dots, k\}$, these are called \textit{transition probabilities}:
\begin{equation}
    p_{ij} = \mathbb{P}(X_{n+1} = j \mid X_n = i) \; ,
\end{equation}
representing the probability of moving from state $i$ to state $j$ in one step. We assume throughout that these probabilities do not depend on the time $n$, a property called \textit{time homogeneity}: the rules of the game do not change as the system evolves. The transition probabilities satisfy, for each state $i$,
\begin{equation}
    p_{ij} \ge 0 \quad \text{for all } j \; , \qquad \sum_{j} p_{ij} = 1 \; ,
\end{equation}
since the system must transition to \textit{some} state at each step, and probabilities must be non-negative and sum to one.

\medskip

These probabilities are naturally arranged into a \textit{transition matrix} $\mathbf{P}$, where the entry in row $i$ and column $j$ is $p_{ij}$. Each row of $\mathbf{P}$ sums to one, which is why such matrices are called \textit{stochastic matrices}. The transition matrix is the complete mathematical description of a Markov chain: given the initial state and $\mathbf{P}$, all probabilistic properties of the chain can in principle be derived. It is a remarkably compact encoding of the system's dynamics.

\medskip

\textit{Example.} Returning to Markov's original setting, consider a two-state chain with states V (vowel) and C (consonant). Suppose that a vowel is followed by another vowel with probability $0.4$ and by a consonant with probability $0.6$, while a consonant is followed by a vowel with probability $0.3$ and by another consonant with probability $0.7$. These numbers are schematic, but they reflect the qualitative structure Markov observed in Russian text: consonants are more likely to follow consonants, and vowels to follow vowels, but neither transition is impossible. The transition matrix is
\[
    \mathbf{P} = \bordermatrix{
        & \text{V} & \text{C} \cr
        \text{V} & 0.4 & 0.6 \cr
        \text{C} & 0.3 & 0.7
    } \; .
\]
Reading row by row: from a vowel, we go to another vowel with probability $0.4$ and to a consonant with probability $0.6$; from a consonant, we go to a vowel with probability $0.3$ and stay a consonant with probability $0.7$. Every row sums to one.

\medskip

\textit{Example.} A simple weather model with two states, Sunny ($S$) and Rainy ($R$), captures the same structure in a different domain. Suppose that a sunny day is followed by another sunny day with probability $0.8$ and a rainy day with probability $0.2$, while a rainy day is followed by a sunny day with probability $0.4$ and another rainy day with probability $0.6$. The transition matrix is
\[
    \mathbf{P} = \bordermatrix{
        & S & R \cr
        S & 0.8 & 0.2 \cr
        R & 0.4 & 0.6
    } \; .
\]
Despite the completely different subject matter, the mathematical structure is identical to the vowel-consonant chain. This is the power of abstraction: once a process has been encoded as a Markov chain, all the tools of the theory apply regardless of whether the states represent letters, weather patterns, or anything else.

\subsection{Chapman--Kolmogorov equations}

A natural question is: what is the probability of moving from state $i$ to state $j$ in exactly $n$ steps, rather than just one? If we want to know the probability that it will be sunny in three days given that it is raining today, we cannot simply read off an entry of $\mathbf{P}$---we need to account for all the possible intermediate states the system might pass through. These are called the \textit{$n$-step transition probabilities}, written $p_{ij}^{(n)}$, and they satisfy the \textit{Chapman--Kolmogorov equations}:
\begin{equation}
    p_{ij}^{(m+n)} = \sum_{k \in \mathcal{S}} p_{ik}^{(m)} \, p_{kj}^{(n)} \; .
\end{equation}
This equation says that the probability of going from $i$ to $j$ in $m+n$ steps equals the sum, over all possible intermediate states $k$, of the probability of reaching $k$ from $i$ in $m$ steps and then reaching $j$ from $k$ in $n$ steps. All paths through the state space are accounted for, weighted by their probabilities. In matrix notation, this takes a particularly clean form: the $n$-step transition matrix is simply $\mathbf{P}^n$, the $n$-th matrix power of $\mathbf{P}$ \cite{norris_1997_markov}. To find the probability distribution over states after $n$ steps, starting from an initial distribution $\boldsymbol{\pi}_0$, one computes $\boldsymbol{\pi}_0 \mathbf{P}^n$.

\subsection{Classification of states}

Not all states in a Markov chain behave in the same way, and understanding the long-run behavior of the chain requires classifying them. Some states, once left, may never be returned to; others are visited again and again, indefinitely. A state $i$ is called \textit{recurrent} if, starting from $i$, the chain returns to $i$ with probability one. Intuitively, a recurrent state is one the system cannot escape permanently---it will always come back, however long it wanders. A state is called \textit{transient} if there is a positive probability that the chain never returns to it after leaving. A particularly extreme case is an \textit{absorbing state}: a state $i$ with $p_{ii} = 1$, from which there is no escape at all. In the gambler's ruin problem, bankruptcy and winning are both absorbing states.

\medskip

A Markov chain is called \textit{irreducible} if every state can be reached from every other state in a finite number of steps. Irreducibility ensures that the chain does not fracture into isolated components that can never communicate. It is a global property of the chain as a whole, and it plays a central role in guaranteeing the existence and uniqueness of a stationary distribution, as we will see shortly. A chain is called \textit{aperiodic} if the system does not get trapped in a fixed cycle, returning to each state only at regular multiples of some period greater than one. Irreducibility and aperiodicity together are the key conditions for the chain to behave well in the long run.

\subsection{Stationary distributions and detailed balance}

One of the most important questions about a Markov chain is what happens as time goes on. If we run the chain for a very long time, does the distribution over states settle into a stable pattern, independent of where we started? The answer is given by the notion of a \textit{stationary distribution}. A probability distribution $\boldsymbol{\pi} = (\pi_1, \pi_2, \dots, \pi_k)$ over the state space is called \textit{stationary} if
\begin{equation}
    \boldsymbol{\pi} = \boldsymbol{\pi} \mathbf{P} \; , \qquad \sum_i \pi_i = 1 \; .
\end{equation}
In words: if the system starts in distribution $\boldsymbol{\pi}$, one step of the chain leaves it in distribution $\boldsymbol{\pi}$. The distribution is unchanged by the dynamics. The stationary distribution represents the long-run proportion of time the chain spends in each state---in the weather model, it tells us what fraction of days, over a very long period, will be sunny and what fraction rainy, regardless of today's weather.

\medskip

\textit{Example.} For the weather chain above, the stationary distribution satisfies $\boldsymbol{\pi}\mathbf{P} = \boldsymbol{\pi}$, which gives the system of equations
\[
    \pi_S = 0.8\,\pi_S + 0.4\,\pi_R \; , \qquad \pi_R = 0.2\,\pi_S + 0.6\,\pi_R \; , \qquad \pi_S + \pi_R = 1 \; .
\]
Solving gives $\pi_S = 2/3$ and $\pi_R = 1/3$. In the long run, two thirds of days will be sunny and one third rainy, regardless of whether today is sunny or rainy.

\medskip

A sufficient condition for $\boldsymbol{\pi}$ to be stationary is the \textit{detailed balance} condition:
\begin{equation}
    \pi_i \, p_{ij} = \pi_j \, p_{ji} \quad \text{for all } i, j \; .
\end{equation}
This condition says that, in the stationary regime, the flow of probability from state $i$ to state $j$ exactly equals the flow from $j$ to $i$. The system is in equilibrium not just globally but at every pair of states individually. Chains satisfying detailed balance are called \textit{reversible}, because their dynamics look the same whether run forwards or backwards in time. Reversibility is a stronger condition than stationarity---it implies stationarity but is not implied by it---and it plays a central role in the design of Markov chain Monte Carlo algorithms.

\subsection{Ergodicity and convergence}

The central theorem of Markov chain theory, which justifies much of what follows in this chapter, concerns the long-run behavior of irreducible and aperiodic chains. Such chains are called \textit{ergodic}, and for them the following holds: regardless of the initial state, the distribution over states converges to the unique stationary distribution as the number of steps grows. Formally,
\begin{equation}
    \lim_{n \to \infty} p_{ij}^{(n)} = \pi_j \quad \text{for all } i, j \; ,
\end{equation}
meaning that after sufficiently many steps, the probability of being in state $j$ approaches $\pi_j$ no matter where the chain started \cite{norris_1997_markov}. This is the Markov chain analogue of the law of large numbers, and it is precisely what Markov needed to answer Nekrasov: even for dependent sequences, long-run averages stabilize, and they stabilize at the values the stationary distribution prescribes. The chain forgets its initial condition, and what remains is the intrinsic geometry of the transition structure.

\medskip

The ergodic theorem also has a time-average interpretation: for any function $g$ of the state, the time average
\[
    \frac{1}{n} \sum_{t=0}^{n-1} g(X_t) \xrightarrow{n \to \infty} \sum_i \pi_i \, g(i) = \mathbb{E}_{\boldsymbol{\pi}}[g] \; ,
\]
with probability one. This is a powerful statement: long-run time averages converge to expectations under the stationary distribution, regardless of starting point. It is this property that makes ergodic Markov chains useful as computational tools---by simulating the chain for long enough, we can estimate expectations under $\boldsymbol{\pi}$ without knowing $\boldsymbol{\pi}$ in closed form. This idea will return in full force in further sections.

\subsection{Examples}

\textbf{The random walk.} Consider a particle on the integers that, at each step, moves one position to the right with probability $p$ and one position to the left with probability $1-p$. This is the simplest example of a random walk, and one of the most studied objects in probability theory. The state space is infinite, but the Markov property holds: the next position depends only on the current one. When $p = 1/2$, the walk is symmetric and recurrent in one and two dimensions---the particle will return to any given position with probability one---but transient in three or more dimensions, where it drifts away and never returns. This celebrated result is due to Pólya \cite{polya_1921_randomwalk}, and it has been described as one of the most surprising facts in all of probability.

\medskip

\textbf{The gambler's ruin.} A gambler starts with $k$ coins and plays a series of fair games, winning or losing one coin at each step. The game ends when the gambler either reaches $N$ coins (winning) or loses everything (ruin). Both endpoints are absorbing states: once reached, the chain stays there forever. The probability of ruin, starting from $k$ coins, is $(N-k)/N$ for a fair game. This result has a sobering implication: even in a perfectly fair game, a gambler with finite resources playing against an opponent with larger resources will eventually go bankrupt with certainty if the game continues indefinitely. The mathematics offers no comfort to the persistent gambler.

\medskip

\textbf{PageRank.} Google's original PageRank algorithm, developed by Brin and Page in the late 1990s, models a web surfer who randomly follows hyperlinks from page to page \cite{brin_1998_pagerank}. The web is treated as a directed graph, and the surfer's position is a Markov chain on the pages of the internet. The stationary distribution of this chain assigns to each page a score proportional to how often the random surfer would visit it in the long run. Pages linked to by many important pages receive high stationary probability, capturing the intuition that influence propagates through networks. The algorithm that shaped the modern internet is, at its mathematical core, the computation of a stationary distribution of a very large Markov chain.

\medskip

\textbf{The Wright--Fisher model.} In population genetics, the Wright--Fisher model describes how the frequency of a gene variant changes from generation to generation due to random reproduction \cite{ewens_2004_genetics}. A population of $N$ individuals is replaced each generation by sampling with replacement from the previous one. The frequency of a variant is the state, and the resulting Markov chain has absorbing states at zero and one, corresponding to the extinction or fixation of the variant. This model provides the mathematical foundation for understanding genetic drift---the random fluctuation of gene frequencies over time---and represents one of the earliest and most influential applications of Markov chains in biology.